\documentclass[12pt,a4paper]{article}
\usepackage[margin=2.5cm]{geometry}
\usepackage{amsmath,graphicx,hyperref,natbib,booktabs,caption,xcolor}
\definecolor{ai4sgreen}{HTML}{0D3B2E}
\definecolor{ai4sgold}{HTML}{C8956C}

\title{\textbf{\textcolor{ai4sgold}{AI4S-Journal} Dataset \& Benchmark}}
\author{Author Name\textsuperscript{1*}, Co-Author\textsuperscript{2} \
\textsuperscript{1}Department, Institution, City, Country \
\textsuperscript{2}Department, Institution, City, Country \
\textsuperscript{*}Email: email@institution.edu}
\date{}

\begin{document}
\maketitle

\begin{abstract}
Describe the new dataset or benchmark, its size, domain, intended use, and
how it advances the field over existing resources.
\end{abstract}
\noindent\textbf{Keywords:} dataset, benchmark, evaluation, [domain]

\section{Introduction}
Motivate the need. Review existing resources and their limitations.

\section{Dataset Description}
\subsection{Data Collection}
How was the data collected/generated? Include ethical considerations.
\subsection{Data Composition}
Size, format, splits (train/val/test), and class distribution.

\begin{table}[h]
\centering
\caption{Dataset Statistics}
\begin{tabular}{@{}lrrr@{}}
\toprule
Split & Samples & Classes & Avg. Length \
\midrule
Training   & 100,000 & 50 & 256 \
Validation &  10,000 & 50 & 256 \
Test       &  10,000 & 50 & 256 \
\bottomrule
\end{tabular}
\end{table}

\subsection{Data Preprocessing}
Document all preprocessing and augmentation steps.

\section{Benchmark Protocol}
\subsection{Task Definition}
\subsection{Evaluation Metrics}
\subsection{Baseline Methods}

\section{Baseline Results}
\section{Discussion and Limitations}
\section{Conclusion}
\section*{Data Availability}
Provide download link and license (CC BY 4.0 recommended).

\bibliographystyle{unsrt}
\begin{thebibliography}{99}
\bibitem{d1} Author, A. (Year). \textit{Title}. Journal.
\end{thebibliography}
\end{document}
